% !TeX program = xelatex
% UTF-8。以 XeLaTeX 连续编译两次即可生成含交叉引用的 PDF。
\documentclass[UTF8,a4paper,11pt,fontset=fandol]{ctexart}
\IfFontExistsTF{SimSun}{\IfFontExistsTF{SimHei}{\setCJKmainfont[BoldFont=SimHei]{SimSun}}{}}{}
\usepackage[top=26mm,bottom=25mm,left=25mm,right=25mm,headheight=14pt,headsep=6mm]{geometry}
\usepackage{amsmath,array,booktabs,tabularx,caption,cite,fancyhdr,hyperref}
\hypersetup{unicode,hidelinks,pdfdisplaydoctitle=true,
  pdftitle={面向 ORCA 密度泛函计算的大语言模型智能体综述},
  pdfsubject={ORCA 智能体的任务规划、工具调用、执行反馈与科学验证},
  pdfkeywords={大语言模型, ORCA, 智能体, 密度泛函理论, 量子化学, 科学工作流, 可复现性}}
\ctexset{section={format=\large\bfseries\raggedright,beforeskip=2.8ex plus .5ex,afterskip=1.0ex},
  subsection={format=\normalsize\bfseries\raggedright,beforeskip=1.8ex plus .4ex,afterskip=.5ex}}
\pagestyle{fancy}\fancyhf{}
\fancyhead[C]{\small 面向 ORCA 密度泛函计算的大语言模型智能体综述}
\fancyfoot[C]{\small\thepage}
\renewcommand{\headrulewidth}{0.3pt}
\setlength{\emergencystretch}{2em}
\setlength{\parskip}{.15em}
\renewcommand{\arraystretch}{1.18}
\setlength{\tabcolsep}{5pt}
\captionsetup[table]{position=top,font=small,labelfont=bf,labelsep=quad,skip=6pt}
\newcolumntype{Y}{>{\raggedright\arraybackslash}X}
\newcolumntype{P}[1]{>{\raggedright\arraybackslash}p{#1}}
\newcommand{\doi}[1]{\href{https://doi.org/#1}{\nolinkurl{#1}}}
\newcommand{\ax}[1]{\href{https://arxiv.org/abs/#1}{arXiv:\,#1}}
\newcommand{\gitrepo}[2]{\href{https://github.com/#1}{#2}}
\newcommand{\snap}[2]{\href{https://github.com/#1/commit/#2}{\texttt{#2}}}
\urlstyle{same}
\raggedbottom
\begin{document}
\thispagestyle{plain}
\begin{center}
  {\LARGE\bfseries 面向 ORCA 密度泛函计算的\\[.3ex]大语言模型智能体综述\par}
\end{center}

\section{引言}
一次分子计算可以正常结束，却仍未回答研究者的问题。
例如，ORCA 能给出目标分子的核磁屏蔽常数；
若要求的是化学位移，工作流还需取得与所用协议一致的参照量\cite{quntur}。
量子化学智能体（agent）需要从自然语言问题中识别目标物理量，
决定需要哪些计算，调用求解器，再根据返回的结果判断原计划是否充分。
大语言模型（large language model，LLM）参与其中哪些决定，
以及软件和研究者如何检查这些决定，是本文比较各项工作的主线。

\section{文献范围与比较方法}
本文考察 17 篇论文或预印本及 1 项独立评测软件。
重点讨论实际以 ORCA 执行密度泛函理论（density functional theory，DFT）计算的案例，
并以 ORCA 输入生成和多参考波函数计算作为相邻证据。
对于支持多种程序的系统，按具体算例采用的求解器与方法归类。
本文比较计划由谁制定、何种反馈会改变下一步，以及哪些证据足以支持报告中的结论。
文中保留对照实验的任务单位、基线和判定依据，仅在同一研究的相同设置内比较数值。

\section{ORCA 智能体如何制定并执行计划}

\subsection{从生成输入到调用求解器}
Jacobs 与 Pollice 的工作为输入生成这一环节提供了基线。
在由真实 ORCA 输入反推的 294 条测试提示上，
GPT-4o 基线生成的输入有 5.872\% 通过执行检查；
采用合成数据微调和思维链提示后，该比例升至 18.027\%。
检查统一附加水分子坐标，并要求 ORCA 在五秒内以零退出码结束。
这个试验回答的是模型能否把请求转换成程序可接受的输入。
原分子的几何、计算之间的依赖以及最终物理量，则要在整项任务中检验\cite{input}。

ArIA 将输入生成接入交互式计算。系统解释用户指令，运行 ORCA，
生成图形和报告，并允许用户针对计算结果继续提问。
作者分别公开了开发和部署材料。
模型开发部分利用 ORCA 试运行以及校验和反馈模块检查候选输入，
部署部分将训练后的模型适配器接入 LangGraph 工作流\cite{aria}。
因此，可以分别评价输入生成的改进和用户通过追问能否正确理解某次计算。
这两种评价各自需要对应的输入、输出与交互记录。

\subsection{Q-Planner 的计划表示与提交前校验}
Q-Planner 将用户请求转换成有向无环图。
模型识别化合物、选择适用的计算流程，并声明每个节点所需的输入和产物。
计划通过结构校验后，执行器按依赖顺序调度工具，
经 ORCA Python Interface 调用 ORCA，并将返回值交给后续节点。
计划可以预先声明错误类型及对应的参数调整；
执行器可并行运行独立任务，并依照这些规则处理已知故障\cite{qplanner}。

氯氟乙酸水相 $\mathrm{p}K_\mathrm{a}$ 的算例揭示了这一分工的作用。
模型第一次规划时漏掉几何优化，校验器在提交 ORCA 前发现缺项。
第二次规划补齐该步骤后，依赖优化结构的自由能计算才得以执行\cite{qplanner}。
在另一项批量规划试验中，对 10 至 100 个分子的单点任务，
模板展开的 30 份计划全部包含所需节点，逐节点生成时则为 28/30。
这里的分母是计划数。它能衡量代理有没有漏掉任务，
逐个分子的能量是否可信仍须依据实际作业及其输出判断。

\subsection{El Agente Q 与 Quntur 的任务委派和运行反馈}
El Agente Q 采用分层委派。上层的计算化学代理将问题拆成步骤，
负责 DFT 的代理再分别将输入编写和集群提交交给专门角色。
各角色调用工具后，把结果或故障返回上层，
使后续步骤能够引用已经生成的文件，并根据失败信息调整安排。
作者公开了六道习题和两项案例的计算档案。
习题已给定泛函和基组，因而主要检验的是代理能否把指定方法落实为一组连贯的 ORCA 作业\cite{agente}。

Quntur 将方法选择和资料检索也纳入这条反馈回路。
编写输入时，负责协调的代理判断当前缺少的是软件语法还是科学依据，
再决定查阅 ORCA 手册、检索文献或委派给领域代理。
读取大型 ORCA 输出时，工具先建立章节及行号索引，
代理据此定位所需表格，抽查表头和数据行，再提取字段用于分析\cite{quntur}。
同一输出可能含多个电子态或多种性质，
章节和行号使研究者可以核查代理选取的数据是否对应所问的物理量。

\section{ORCA 输出如何改变智能体的科学判断}
Quntur 对 4-二甲氨基苯甲腈（DMABN）的分析提供了一个实例。
代理用优化几何和自然跃迁轨道判断两个低能单重态的局域激发或电荷转移特征。
发现所得态特征的变化与文献描述不符后，代理再次检索文献，
追查到本次计算采用气相，而被比较的现象涉及溶液条件\cite{quntur}。
这一步反馈改变了报告可以作出的解释，也指明进一步计算需要控制的环境条件。

在 17 道任务各独立重复五次的对照中，Quntur 与启用 ORCA 技能的 Claude Code
均使用 Claude Opus 4.5。专家按规划、几何、输入与执行、后分析评分，
两组平均分别为 90.5 和 61.7 分（满分 100）。
在核磁与磁性任务中，对照组五次中有四次遗漏了一项或两项参照计算\cite{quntur}。
这一实例把评分落实到一项可检查的代理决定，即是否为所求性质安排了必要的参照。
整体均分反映完整配置的表现，各组成模块的独立贡献仍需消融实验估计。

Lee 与 Rondinelli 的 ORCA 多参考波函数研究提供了另一种评测设计。
代理负责选择活性空间、提交作业、处理报错并匹配目标电子态。
在 558 个参考跃迁中，加入专家决策指导后，匹配数由 139 增至 246；
在双方共同匹配的 109 项上，平均绝对误差则由 0.242 变为 0.268~eV\cite{mr}。
这项相邻方法研究说明，评价智能体应同时报告取得可用结果的目标数，以及同一批已匹配目标上的数值误差。

\section{其他计算平台对智能体设计的启示}
多后端分子计算系统提供了工具分工的对照。
ChemGraph 将任务拆解后交由 ASE 工具执行，小分子几何优化基准的参照采用 NWChem。
Aitomia 通过 MLatom 选择方法并执行脚本，
其己醇红外示例采用 AIQM2\cite{chemgraph,aitomia}。
Cloudra 用反应路径端点检验 Gaussian 找到的鞍点；
VirtualLab\_CC 则由项目负责人审批是否进入下一计算阶段。
后者的公开方法协议将构象 DFT 筛选交给 Gaussian，将 DLPNO-CCSD(T) 参照计算交给 ORCA\cite{cloudra,virtuallab}。
这些案例分别说明了工具选择、独立验证和人工复核如何参与同一计算任务。

周期体系研究集中说明代理应在什么时刻再次介入。
VASPilot 根据作业日志修订输入，LatticeMind 在提交前检查步骤间的数据依赖\cite{vaspilot,lattice}。
AutoDFT 的同模型对照中，闭环反馈使 34 项 VASP 任务的完成数由 28 增至 32，
每项平均 H100 GPU 小时由 22.1 增至 44.0\cite{autodft}。
CatMaster 依据 DFT 复核结果调整候选范围，MOFGen 将生成、筛选和复核分阶段组织\cite{catmaster,mofgen}。
在评测方面，DFT-Agent Bench 和 CompChemBench 分别检查求解器执行记录与工作流审计能力。
泛函搜索研究则用留出数据检验代理所提候选的预测能力\cite{dftbench,compbench,xc}。

\section{讨论与局限}
上述 ORCA 工作提示，反馈的价值取决于它能纠正哪一类决定。
程序试运行可以暴露输入语法错误，计划校验可以发现上游几何优化的缺失。
识别参照态遗漏或气相与溶液条件混用，则需要把实际输出重新对照用户的问题。
对于依赖关系已知的任务，Q-Planner 的对照支持在计划生成时展开模板；
对于解释取决于中间结果的任务，Quntur 的案例给出了重新查阅资料的具体理由\cite{qplanner,quntur}。

现有研究采用的题目、专家量规和计算预算不同，
单个系统的高分难以直接外推为另一类任务上的优势。
更有辨别力的比较，应在同一模型、工具和预算下记录代理作出的关键选择，
同时报告任务覆盖、共同样本上的性质误差、重复计算和资源消耗。
这样才能判断收益来自更完整的计划、更有针对性的修复，还是更多的计算机会\cite{quntur,mr,autodft}。

\section{项目后续发展方向}
对于基于 ORCA 的分子计算代理，研究者在一次作业结束后往往面临两个连续的判断。
已有结果能回答哪些后续问题？
如果确需补算，有限的机时应优先分配给哪些候选？
下文以项目存档中的乙醇和正己烷计算为例提出两项研究设计。
评估所需的盲测样本与对照实验均有待建立。

\subsection{面向后续提问的既有计算证据复用}
当研究者改变所问的物理量，同一组输出能够支持的结论也会改变。%
Quntur 的核磁共振评测表明，即使已获得目标分子的屏蔽常数，
遗漏参照计算仍会影响化学位移的推导\cite{quntur}。%
Lo Scudo 在热化学工作流中进一步检验了物种记录的完整性，
在缺少必要物种的热化学数据时阻止了反应量的汇总\cite{tracestate}。
后续研究拟检验，当原始文件保持不变而研究者改变提问时，
代理能否准确判断证据的适用范围，并指出推导新结论所缺的条件或数据。

项目的乙醇算例说明了这一问题的实际含义。
在相同的分子几何上，以 B3LYP 分别计算气相与 CPCM 水环境中的单点电子能。
按 $\Delta E=E_{\mathrm{water}}-E_{\mathrm{gas}}$ 定义，%
def2-SVP 和 def2-TZVP 基组给出的差值分别为 $-16.61$ 和 $-18.28$~kJ\,mol$^{-1}$。
因此，在所考察的几何、泛函与两种基组下，水环境中的电子能均较低，
两组差值相差约 $1.67$~kJ\,mol$^{-1}$\cite{projectarchive}。
若提问改为标准水合自由能，则须明确温度、标准态及溶剂化自由能的计算方案，
并核对已有输出覆盖了哪些贡献。

拟以未参与开发的分子构造成对盲测。
每对问题共用一组封存的 ORCA 输出，只改变差值方向、参照条件或目标物理量。
由独立标注者给出可从原件直接读出或经后处理求得的答案，
并标明哪些问题需要补充条件或启动新计算。
在相同模型、工具和预算下，将现有系统与逐项核对目标量、计算条件和数值出处的方案比较。
主要指标为两问的判断与物理解释均正确的题对比例，
同时报告证据不足却给出数值的比例、证据充分却未作答的比例，以及可避免的 ORCA 启动次数。

\subsection{有限预算下的自适应构象选择}
在多层级构象筛选中，较便宜的方法可为大量候选排序，再以 DFT 精化其中一部分。%
CENSO 已将这一策略用于柔性分子的构象系综与自由能计算\cite{censo}。%
VirtualLab\_CC 记录了构象筛选与后续精化的机时，AutoDFT 的对照实验也表明，
允许代理依据中间结果修订方案会改变计算成本\cite{virtuallab,autodft}。
后续研究拟检验，在高层级单点计算次数相同的条件下，
利用已返回的能量调整下一次选择，能否提高识别低能构象的概率。

项目已完成两批正己烷气相 $r^2$SCAN-3c 单点计算，
分别获得四个和六个固定几何的批内电子能排序\cite{projectarchive}。
拟另选小分子，每个分子固定多于四个候选几何，在统一参考层级算全单点能并封存。
各策略均可见两种低成本方法的排序，并有四次查询参考能量的机会。
基线一选取主筛选方法排序最低的四个候选，
基线二按预定义的二面角类别兼顾低能候选的构象家族覆盖。
自适应策略先选主筛选方法的最低能候选，
再依据已返回的能量、两种排序出现分歧的位置和尚未覆盖的家族选择下一项。

揭盲后，按分子统计参考最低能构象的命中率。
设 $S$ 为固定候选集，$\hat c$ 为最终选定的构象，
能量差 $E_{\mathrm{ref}}(\hat c)-\min_{c\in S}E_{\mathrm{ref}}(c)$ 衡量漏选最低能构象的代价。
并列报告实际 CPU 核时与模型词元用量，再按两种低成本排序的一致程度分组，
可检验自适应选择在哪些分子上有收益。
结论限于给定几何及参考层级下的电子能排序。

\section{结论}
ORCA 工作流中已有两类经实验证实的智能体设计。
依赖关系明确时，模型可以在提交前形成可校验的计划，
由确定性执行器展开批量任务并处理预先定义的错误。
科学问题随中间结果变化时，代理需要保留原始输出，
追查缺失的参照或条件，再决定是否补算以及如何解释结果。
相应的比较应记录计划在哪一步被纠正、哪些判断改变了后续行动，
以及这些决定对科学结论和计算成本的实际影响\cite{qplanner,agente,quntur}。

\begingroup
\small
\begin{thebibliography}{99}
\raggedright
\setlength{\itemsep}{.35em plus .1em}
\bibitem{quntur}
  PÉREZ-SÁNCHEZ J B, ZOU Y, CAMPOS-GONZALEZ-ANGULO J A, et al.
  El Agente Quntur: A research collaborator agent for quantum chemistry [J/OL].
  \emph{Chem}, 2026-09-24（在线发表）, 103259.
  DOI: \doi{10.1016/j.chempr.2026.103259}.
\bibitem{input}
  JACOBS P F, POLLICE R.
  Developing large language models for quantum chemistry simulation input generation [J].
  \emph{Digital Discovery}, 2025, 4(3): 762--775.
  DOI: \doi{10.1039/D4DD00366G}.
\bibitem{aria}
  CHANMUNGKALAKUL S, VAN DER REE M, GIUNTOLI A, et al.
  Molecular Simulations Assisted by an Artificial Intelligence Agent (ArIA) [EB/OL].
  \emph{ChemRxiv}, v2, 2026-05-08.
  DOI: \doi{10.26434/chemrxiv.15002344/v2}.
\bibitem{qplanner}
  ZHANG R, GAO W, KURASHIGE Y.
  Q-Planner: A Harness System for Quantum Chemistry Agents [J/OL].
  \emph{ACS Omega}, 2026-11-10（在线发表）.
  DOI: \doi{10.1021/acsomega.6c05744}.
\bibitem{agente}
  ZOU Y, CHENG A H, ALDOSSARY A, et al.
  El Agente: An autonomous agent for quantum chemistry [J].
  \emph{Matter}, 2025, 8(7): 102263.
  DOI: \doi{10.1016/j.matt.2025.102263}.
\bibitem{mr}
  LEE V C, RONDINELLI J M.
  Can Autonomous LLM Agents Execute Multireference Quantum Chemistry Calculations?
  [EB/OL].
  \emph{arXiv}, v1, 2026-09-11.
  \ax{2609.13357v1}.
\bibitem{chemgraph}
  PHAM T D, TANIKANTI A, KEÇELİ M.
  ChemGraph as an agentic framework for computational chemistry workflows [J].
  \emph{Communications Chemistry}, 2026, 9(1): 33.
  DOI: \doi{10.1038/s42004-025-01776-9}.
\bibitem{aitomia}
  HU J, NAWAZ H, HOU Y-F, et al.
  Aitomia: An Agentic Framework for AI-Driven Atomistic and Quantum Chemical Simulations
  [J].
  \emph{Journal of Chemical Theory and Computation}, 2026, 22(13): 6557--6571.
  DOI: \doi{10.1021/acs.jctc.6c00591}.
\bibitem{cloudra}
  YANG C, ZHOU Q, LI Y.
  Cloudra: Geometric Reasoning and Scientific Validation for Autonomous Reaction Mechanism
  Exploration [EB/OL].
  \emph{arXiv}, v1, 2026-11-02.
  \ax{2611.02909v1}.
\bibitem{virtuallab}
  KIENINGER M, SEGOVIA M, VEGA-TEIJIDO M, et al.
  VirtualLab\_CC: An LLM-augmented virtual laboratory for automated computational
  chemistry research [J].
  \emph{Artificial Intelligence Chemistry}, 2026, 4(2): 100129.
  DOI: \doi{10.1016/j.aichem.2026.100129}.
\bibitem{vaspilot}
  LIU J, ZHU T, YE C, et al.
  VASPilot: MCP-facilitated multi-agent intelligence for autonomous VASP simulations [J].
  \emph{Chinese Physics B}, 2025, 34(11): 117106.
  DOI: \doi{10.1088/1674-1056/ae0681}.
\bibitem{lattice}
  VEGH N M, ZHU Y, PRENTKI R J, et al.
  LatticeMind: A Validation-Gated Agentic Framework for Density Functional Theory
  Workflows [EB/OL].
  \emph{arXiv}, v1, 2026-07-19.
  \ax{2610.26869v1}.
\bibitem{autodft}
  YANG P, ZHANG Z, LI Y, et al.
  AutoDFT: A Closed-Loop Multi-Agent Framework for Autonomous DFT Calculations [EB/OL].
  \emph{arXiv}, v2, 2026-06-04.
  \ax{2605.26179v2}.
  \href{https://phyang.top/}{KDD 2027 已接收}。
\bibitem{catmaster}
  CHEN H, QIU J, TEW Y S, et al.
  Autonomous computational catalysis through an agentic research system [EB/OL].
  \emph{arXiv}, v4, 2026-06-05.
  \ax{2601.13508v4}.
\bibitem{mofgen}
  INIZAN T J, YANG S, KAPLAN A, et al.
  System of agentic AI for the discovery of metal--organic frameworks [J].
  \emph{Nature Computational Science}, 2026, 6(11): 1223--1234.
  DOI: \doi{10.1038/s43588-026-01058-x}.
\bibitem{dftbench}
  LIU S, YANG J, WAN W, et al.
  DFT-Agent Bench: Scaffold Effects and Execution Provenance in Autonomous DFT [EB/OL].
  \emph{arXiv}, v1, 2026-11-07.
  \ax{2611.10629v1}.
\bibitem{compbench}
  WU C.
  CompChemBench: A Benchmark for AI Agents on Computational Chemistry Tasks [CP/OL].
  v1.1. \emph{GitHub}, 2026 [2026-11-22].
  \href{https://github.com/Chenghao-Wu/compchem-bench/tree/772d3f3e1a93}{版本快照 772d3f3e1a93}.
\bibitem{xc}
  CHANDRASEKHAR A, BARATI FARIMANI A.
  Agent-Guided Discovery of Exchange-Correlation Functionals with Physics-based Validation
  [EB/OL].
  \emph{arXiv}, v1, 2026-11-09.
  \ax{2611.13329v1}.
\bibitem{tracestate}
  LO SCUDO F.
  Trace-State Validation for Agentic Computational Chemistry Workflows [EB/OL].
  \emph{ChemRxiv}, v1, 2026-07-20.
  DOI: \doi{10.26434/chemrxiv.15006307/v1}.
\bibitem{projectarchive}
  BG6022-V4 项目组.
  正己烷构象与乙醇环境效应计算记录及报告复核 [R/OL].
  \emph{GitHub}, 2026.
  \href{https://github.com/az87988799/BG6022-v4/blob/1e8795b0b10deb8f288eccb9d95befb742958ff9/FinalReport-LaTeX/sections/06a-calculation-results.tex}
    {计算结果档案}，%
  \href{https://github.com/az87988799/BG6022-v4/blob/1eefccdb742ed81692900f9660a8db950ae1f273/docs/acceptance/live-failure-repair-20261021/repair-report.md}
    {乙醇运行复核记录}。
\bibitem{censo}
  GRIMME S, BOHLE F, HANSEN A, et al.
  Efficient Quantum Chemical Calculation of Structure Ensembles and Free Energies for
  Nonrigid Molecules [J].
  \emph{The Journal of Physical Chemistry A}, 2021, 125(19): 4039--4054.
  DOI: \doi{10.1021/acs.jpca.1c00971}.
\end{thebibliography}
\endgroup

\clearpage
\appendix
\setcounter{table}{0}
\renewcommand{\thetable}{\Alph{section}.\arabic{table}}
\renewcommand{\theHtable}{appendix.\Alph{section}.\arabic{table}}
\section{公开软件与复现材料}
\label{app:repos}
ORCA 智能体公开的复现材料各有侧重。
Q-Planner 的代码快照包含技能说明、计划校验器、任务调度器和 ORCA 工具服务，
可对照计划文件检查节点的输入与产物\cite{qplanner}。
\href{https://doi.org/10.5683/SP3/JU2BQK}
  {El Agente Q 的八项任务档案}和%
\href{https://doi.org/10.5683/SP3/RDSOEA}
  {Quntur 与对照系统的评测档案}
保存运行产物，适合复查代理的决策与最后的数值是否对应。
Quntur 作者另行说明其代理源码未公开\cite{agente,quntur}。
\href{https://doi.org/10.34894/RNPTDS}
  {ArIA 的部署与训练材料}可用来追踪输入生成和反馈模块。
Jacobs 与 Pollice 提供%
\href{https://git.lwp.rug.nl/pollice-research-group/ORCAInputFileSynthesis/-/blob/f60237db580d44849f30179deed090317a4a768a/Classes/Helpers/ORCARunner.py}
  {ORCA 测试程序}和%
\href{https://doi.org/10.34894/WNRHA4}{数据归档}\cite{aria,input}。

表~\ref{tab:repos} 汇总 12 个公开 GitHub 仓库。
Q-Planner 使用论文标签 \texttt{qplanner-acsomega-r2.2}；
各条目的提交标识链接到检索时核对的版本。

\begin{table}[htbp]
\centering
\caption{纳入比较的公开代码与数据仓库（检索于 2026 年 11 月 22 日）}
\label{tab:repos}
\small
\begin{tabularx}{\linewidth}{@{}P{3.45cm}YP{2.95cm}@{}}
\toprule
软件或数据集
  & 可复查的代理或评测环节
  & 提交标识\\
\midrule
\gitrepo{zenqiaz/Q-planner}{Q-Planner}
  & 计划图、校验器与 ORCA 工具调用日志\cite{qplanner}
  & \snap{zenqiaz/Q-planner}{570409b96fa8}\\
\gitrepo{argonne-lcf/ChemGraph}{ChemGraph}
  & 任务分解、ASE 工具调用与结果汇总\cite{chemgraph}
  & \snap{argonne-lcf/ChemGraph}{dcb129ef6269}\\
\gitrepo{dralgroup/aitomia}{Aitomia}
  & 方法选择、脚本生成与 MLatom 执行\cite{aitomia}
  & \snap{dralgroup/aitomia}{e9d62acafa53}\\
\gitrepo{yyywaa/Cloudra-mechanism-agent}{Cloudra}
  & 机理任务规划与路径端点验证\cite{cloudra}
  & \snap{yyywaa/Cloudra-mechanism-agent}{d76d5cca194c}\\
\gitrepo{CCBG-Lab/VirtualLab_CC}{VirtualLab\_CC}
  & 研究角色、人工审核与%
    \href{https://github.com/CCBG-Lab/VirtualLab_CC/blob/1c1a58cc8df6/framework/methods/METH_002_criegee_terpene_pes.md}
    {计算协议}\cite{virtuallab}
  & \snap{CCBG-Lab/VirtualLab_CC}{1c1a58cc8df6}\\
\gitrepo{JiaxuanLiu-Arsko/VASPilot}{VASPilot}
  & 作业监控与基于日志的错误修复\cite{vaspilot}
  & \snap{JiaxuanLiu-Arsko/VASPilot}{2d9b2dfaf6e0}\\
\gitrepo{penghui-yang/AutoDFT}{AutoDFT}
  & 闭环执行与开环对照\cite{autodft}
  & \snap{penghui-yang/AutoDFT}{4c8e16ff59b7}\\
\gitrepo{nathaniel-vegh/latticemind-data-package}{LatticeMind 数据包}
  & 输入依赖校验与消融数据\cite{lattice}
  & \snap{nathaniel-vegh/latticemind-data-package}{2a85c9ef5c2b}\\
\gitrepo{q734738781/CatMaster}{CatMaster}
  & 催化剂候选筛选与复核工具\cite{catmaster}
  & \snap{q734738781/CatMaster}{06b856814f98}\\
\gitrepo{tjaffrel/mofgen}{MOFGen}
  & 候选生成与分阶段筛选\cite{mofgen}
  & \snap{tjaffrel/mofgen}{841c93995a4c}\\
\gitrepo{Chenghao-Wu/compchem-bench}{CompChemBench}
  & 任务执行与科学审计评分\cite{compbench}
  & \snap{Chenghao-Wu/compchem-bench}{772d3f3e1a93}\\
\gitrepo{BaratiLab/DFT_functional_agentic_searcher}{交换相关泛函搜索}
  & 候选泛函提出与独立数据验证\cite{xc}
  & \snap{BaratiLab/DFT_functional_agentic_searcher}{741bd205c5a9}\\
\bottomrule
\end{tabularx}
\end{table}

\end{document}
